Retention of micronutrients in fortified foods is highly dependent on the different storage conditions. The strength of this effect is consequently determined by preceding factors like fortification method and the food matrice. This phenomena is clearly to be seen in studies from \citet{chepkoech}, \citet{noodle} or \citet{gari} where they investigated the connection between pre- and post-fortification. Being aware of this, basic research of \citet{li2coatsalt} and \citet{li3capsalt} contributed to a deeper understanding of the underlying mechanisms and consequently, studies from \citet{chepkoech} and \citet{sharma} contributed data on the different storage conditions. Ultimately, they show that there are strong implications on interacting effects resulting in changes for micronutrient retention. Thus, there is a high demand on further research to make fortification strategies more successful. 